\input{./._relpath-to-latexroot.ltx}
\documentclass[\latexroot/\projectname]{subfiles}
\input{\latexroot/@resources/subfile-setup.ltx}
\begin{document}

\section{Conclusion}
\whenintegrated{\label{sec:conclusion}}

This paper studies the macroeconomic effects of both conventional (target-rate) and unconventional (forward-guidance, FG) monetary policy innovations in Mexico using high-frequency surprises around policy announcements. A central message is that raw interest rate surprises need not be clean measures of exogenous MP shocks: they can embed (i) CBI revealed about fundamentals and (ii) systematic RN components that reflect an imperfectly understood reaction function. When these informational components are not separated, contractionary target-rate shocks exhibit the familiar output, price, and exchange-rate puzzles and are followed by a counterintuitive rise in inflation expectations, while contractionary FG shocks appear to generate disproportionately large and fast real effects.

To address this issue, we combine projections on pre-announcement macro-financial information with high-frequency co-movements in equity prices to disentangle pure monetary policy innovations from RN and CBI shocks. Once CBI and RN components are isolated, the estimated effects of contractionary target-rate shocks become broadly consistent with standard transmission mechanisms, featuring delayed contractions in output and prices, an appreciation of the domestic currency, and well-anchored inflation expectations. For FG, isolating RN shocks substantially attenuates the apparent “forward guidance puzzle”: the implied macroeconomic effects become more conservative and broadly similar to those of pure target-rate shocks.

A key contribution is that RN components play \emph{asymmetric} contamination roles across policy instruments. In our estimates, RN shocks embedded in target-rate surprises tend to attenuate the contractionary effects of conventional MP shocks, while RN shocks embedded in FG surprises instead amplify and accelerate the effects of FG innovations. Beyond distinguishing conventional from unconventional tools, the results also speak to the exchange-rate puzzle in EMEs by highlighting how omitted informational components can distort estimated exchange-rate responses to policy tightenings.

Overall, the findings underscore that credible measurement of MP shocks in EMEs requires explicitly accounting for informational and systematic policy-response components in high-frequency surprises. Doing so yields impulse responses that are more interpretable, more comparable across policy tools, and more consistent with conventional monetary transmission.

\smartbib

\pdfonly{
  \captionsetup[figure]{list=no}
  \captionsetup[table]{list=no}
}

\end{document}